\documentclass[12pt]{article}
\usepackage[utf8]{inputenc}
\usepackage[T1]{fontenc}
\usepackage[margin=1in]{geometry}
\usepackage{setspace}
\usepackage{xcolor}
\usepackage[colorlinks=true, urlcolor=blue!60!black, linkcolor=blue!60!black]{hyperref}
\setlength{\parindent}{0em}
\setlength{\parskip}{0.6em}
\renewcommand{\baselinestretch}{1.0}
\begin{document}
\begin{center}
{\LARGE \textit{United States v. Curtiss-Wright Export Corp.}, 299 U.S. 304 (1936)\par}
\vspace{0.5em}
{\large Case Brief --- Two-Page Edition\par}
\end{center}

\section{Memory Jogger}
Congress may give the President far broader discretion in foreign affairs than at home; the President is the ``sole organ'' of external relations (\href{https://supreme.justia.com/cases/federal/us/299/304/}{299 U.S. 304}).

\section{Detailed Case Facts}
A 1934 Joint Resolution made it a crime to sell arms to Bolivia or Paraguay, then at war in the Chaco, if the President found an embargo might help restore peace and so proclaimed. Roosevelt did, and Curtiss-Wright Export Corp. was indicted for conspiring to sell 15 machine guns to Bolivia.

\section{Procedural History}
The district court sustained a demurrer, holding the Resolution an invalid delegation under \textit{Panama Refining} and \textit{Schechter}. The United States appealed directly.

\section{Judicial Votes}
Reversed 7--1. Sutherland wrote for the Court; McReynolds dissented without opinion; Stone did not participate.

\section{Holding}
The Resolution was a valid delegation because external affairs permit broader presidential discretion. The proclamation was sufficient, and its revocation did not extinguish liability.

\section{Analysis of Opinions}
\textbf{Majority:} Enumerated powers are true ``only in respect of our internal affairs''; external sovereignty passed from the Crown to the United States collectively. The President, as ``the sole organ of the nation in its external relations,'' has the information, secrecy, and dispatch foreign affairs require, and embargo statutes since 1794 confirm the practice. \textbf{Later:} Jackson's \textit{Youngstown} concurrence and \textit{Zivotofsky v. Kerry} (2015) confined the ``sole organ'' language.

\section{Examples: Future Applications}
\textbf{Same-Side:} a presidential ban on chip exports upon a strategic-threat finding. \textbf{Opposite-Side:} a presidential order binding state courts without a statute (\textit{Medell\'in}). \textbf{Fence-Sitter:} banning a foreign-owned app used by millions of Americans.

\section{Critique}
Charles Lofgren and others show the States exercised external powers before 1789. Marshall's ``sole organ'' meant communicator, not policymaker, and federal power over foreign affairs does not prove presidential power.

\section{Key Quotations}
\begin{quote}
``It is important to bear in mind that we are here dealing not alone with an authority vested in the President by an exertion of legislative power, but with such an authority plus the very delicate, plenary and exclusive power of the President as the sole organ of the federal government in the field of international relations---a power which does not require as a basis for its exercise an act of Congress but which, of course, like every other governmental power, must be exercised in subordination to the applicable provisions of the Constitution.''
\end{quote}
\noindent --- Justice Sutherland, 299 U.S. at 319--320.

\end{document}
